\chapter{Biological Aging and Repair}
\label{ch:27}

The sixth part of the book takes up the proposal that most sharply distinguishes the program from the schemes of the preceding parts, which is that human lives might be extended far beyond their present span, in the most extreme version to astronomical durations. The proposal has two halves, a biological one and a reliability one. The biological half asks whether the processes that accumulate damage in living tissue can be slowed, arrested, or reversed, and by how much. The reliability half asks what would follow if they could, since a person who does not age is nonetheless exposed to accident, violence, infection, and catastrophe. This chapter treats the first, the next the question of what it would mean for a person to continue, and Chapter 29 the second. The order matters because the biological half is the only one on which there is an experimental literature, and a book that has been careful to grade the readiness of its engineering programs should apply the same discipline here. The state of the evidence is that the biology of aging is a serious and growing science, that interventions exist which extend life in laboratory animals, that none has been shown to extend human lifespan, and that the leap from modest extensions to the indefinite lives of the more ambitious proposals is a leap of extrapolation and not of observation. The evolutionary account of why organisms age, in which maintenance is traded against reproduction, is the disposable soma theory~\cite{kirkwood1977}, and it frames the question of whether repair is in principle available.

Aging, in the biological sense, is the progressive loss of function and the rising probability of death that accompanies the accumulation of damage and the failure of maintenance. The review of López-Otín and colleagues~\cite{lopezotin2013,lopezotin2023}, published in 2013 and revised in 2023, organized the field around a list of hallmarks: genomic instability, telomere attrition, epigenetic alterations, loss of proteostasis, deregulated nutrient sensing, mitochondrial dysfunction, cellular senescence, stem cell exhaustion, and altered intercellular communication, to which later versions added disabled macroautophagy, chronic inflammation, and dysbiosis. The list is a map and not a theory. The hallmarks interact, their relative importance is disputed, and for most it remains unclear whether the hallmark is a cause of aging, a consequence, or a correlate. What the formulation has done is to turn a vague phenomenon into a set of specific processes, each of which can be studied and in principle modified, and the research on each is a research on a well-posed problem in the sense of Chapter 4: a target, a measurable acceptance condition, and a cost of evaluation, which in this case is very high because the relevant outcome is lifespan or healthspan, measured over years or decades.

The interventions with the strongest laboratory evidence fall into a few classes. Restriction of caloric intake, without malnutrition, extends life in yeast, worms, flies, and rodents and appears to improve some markers of health in primates, though the effect on primate lifespan has been disputed between the two long studies conducted. The drug rapamycin, which inhibits a nutrient-sensing pathway, extends lifespan in mice in several independent laboratories, and the finding has been replicated at multiple sites in a protocol designed for the purpose. The removal of senescent cells, which accumulate with age and secrete inflammatory factors, extends healthy lifespan in mice by genetic and by pharmacological means, and drugs of this class, called senolytics, have entered early clinical testing. Reprogramming of cells toward a more youthful epigenetic state by transient expression of certain factors has reversed some markers of age in cells and in animals, with the risk, evident in the same experiments, of loss of cellular identity and of cancer. And the study of species with unusual longevity or with negligible measured senescence, such as some rockfish, the naked mole rat, and certain bats, has identified candidate mechanisms of resistance to cancer and of tissue maintenance. The effects in the best-studied model, the mouse, are of the order of a tenth to a third in extension of lifespan, with effects often larger in worms and flies, smaller or absent in primates, and unproven in humans.

It is important to be clear about what such results do and do not support. They support the proposition that the rate of aging is not fixed by physical law and that it can be modified, within a species, by genetic and pharmacological means. They do not support the proposition that aging can be eliminated, nor that the modest effects seen in short-lived animals will scale to a species whose natural lifespan is eight or nine decades and for which the interaction among causes of death is more complex. The longest documented human lifespan, that of Jeanne Calment, who died in 1997 at the age of 122, has not been exceeded in the intervening decades, and demographic analyses suggest a plateau in the maximum age at death in spite of large reductions in mortality at younger ages. There is a debate about whether the plateau reflects a biological limit or the paucity of data from the very old, and the program takes no position in it, observing only that a claim of the possibility of lives of centuries or more is a claim of the prospective success of research not yet done.

A feature of the research enterprise bears directly on the principle of the book. The evaluation of an intervention against aging is hampered by the lack of an accepted surrogate endpoint. In the regulation of medicines, a drug is approved for a disease, with an endpoint that can be measured in a trial of reasonable length, and aging is not classified as a disease. The result is that trials in humans must use endpoints of the diseases of aging, or biomarkers whose relationship to mortality is imperfectly established, among them epigenetic clocks, which estimate biological age from patterns of chemical marks on the genome. An epigenetic clock that is shifted by an intervention provides evidence that something changed, but not that the change corresponds to a longer or healthier life, and the risk of treating such a surrogate as the goal is the standard one of the substitution of a measure for the thing measured. The role proposed for capable systems in this domain is the one described in the second part. They may help to search the large space of candidate interventions, to model the networks of interacting pathways, and to design trials. They cannot shorten the interval over which lifespan is observed, and the verification of a claim of extension requires the participation of human subjects over long periods, under consent and oversight, in studies whose designers have no financial interest in the outcome.

The ethical and institutional questions are serious and are taken up in Chapter 30. Here it suffices to note two that bear on the biology itself. The first is that research on aging is likely to be pursued by private interests with an incentive to claim benefits before they are demonstrated, in a market for products whose claims are weakly regulated, and the history of anti-aging products is a history of exaggeration and fraud. The standard of the book, independent verification of claims before reliance on them, applies here with full force. The second is the question of distribution. A therapy that extends life by years would, if it were costly, enlarge the difference in life expectancy that already separates rich from poor in most countries, and the principle that the admissible range should be assessed by the persons affected requires that distribution be considered at the outset and not after the technology exists.

\section*{Formalization}

The standard quantitative framework is the Gompertz--Makeham law of mortality~\cite{gompertz1825,makeham1860}, in which the hazard rate at age $x$ is the sum of an age-independent term and an exponentially increasing term, $\mu(x)=A+Be^{\gamma x}$, with $A,B,\gamma>0$. The first term represents extrinsic causes of death (accidents, violence, some infections); the second represents senescence, with a mortality-rate doubling time of $\ln2/\gamma$.

\begin{proposition}
The survival function is $S(x)=\exp\!\big[-Ax-(B/\gamma)(e^{\gamma x}-1)\big]$. (i) If senescence is removed ($B=0$) the survival function is $S(x)=e^{-Ax}$, with median $\ln2/A$ and mean $1/A$, so that lifespan is bounded in expectation by the extrinsic hazard. (ii) If aging is slowed by a factor $k>1$ in the sense $\gamma\mapsto\gamma/k$ with the hazard at a reference age $x_0$ held fixed, the senescent term at age $x$ becomes $B_ke^{\gamma x/k}$ with $B_k=Be^{\gamma x_0(1-1/k)}$, and the two senescent terms coincide at $x_0$, and for $x>x_0$ the modified term is strictly smaller than $Be^{\gamma x}$. (iii) The median age at death $m$ solves $Am+(B/\gamma)(e^{\gamma m}-1)=\ln2$, and for $B\gamma^{-1}e^{\gamma m}\gg Am$ it satisfies $m\approx\gamma^{-1}\ln[\gamma\ln2/B+1]$.
\end{proposition}

\begin{proof}
Integrate the hazard: $S(x)=\exp(-\int_0^x\mu)$, which gives the formula; the case $B=0$ is immediate. For (ii), the new term equals the old at $x_0$ by construction, and $B_ke^{\gamma x/k}=Be^{\gamma x_0}e^{(\gamma/k)(x-x_0)}<Be^{\gamma x_0}e^{\gamma(x-x_0)}=Be^{\gamma x}$ for $x>x_0$. For (iii), set $S(m)=\tfrac12$ and neglect the term $Am$.
\end{proof}

A numerical illustration with parameters chosen to reproduce the mortality of a high-income population: a mortality-rate doubling time of eight years ($\gamma=0.0866\ \mathrm{yr^{-1}}$), an extrinsic hazard $A=5\times10^{-4}\ \mathrm{yr^{-1}}$, and a total hazard of $10^{-3}$ at age 30, whence $B=3.72\times10^{-5}$. The survival to ages $50,70,80,90,100,110$ is $0.944,0.803,0.620,0.336,0.079,0.003$, the median age at death $84.5$ years and the mean $81.1$ years. Slowing aging by a factor $k=1.5$ (doubling time $12$ years) yields a median of $104.6$ years and a mean of $99.7$; slowing by $k=2$ (doubling time $16$ years) yields a median of $122.5$ years and a mean of $116.4$. Removing senescence entirely leaves a median of $\ln2/A=1386$ years and a mean of $2000$ years: the extrinsic hazard sets the ceiling. To obtain a median of a million years would require $A\le6.9\times10^{-7}\ \mathrm{yr^{-1}}$, a figure that is the subject of Chapter 29. The illustration shows both what the biology could achieve if its most optimistic prospects were realized, which is a large change in the span of life but within the range of hundreds of years, and that the larger numbers can be reached only if the extrinsic hazard is reduced by orders of magnitude, which is a problem of environment, institutions, and engineering, and not of biology.
